\scanpage{67}
\mspage{21}
\[
(23)\qquad I\xrightarrow{\underline{d}}\mathbb Z
\qquad (\text{i.e } I=\coprod_m I_m)
\]
jouant le rôle d'une fonction dimension
(on aurait une présentation en quelque sorte
duale avec une fonction type codimension),
satisfaisant la condition
\[
(24)\qquad\text{Si }i,j\in I\,\deleted{\unsure{\varnothing}},\ i\leq j,\text{ alors }d(j)=d(i)+1\Longleftrightarrow j\text{ succ. de }i
\]
Posons alors, pour $m,n_{\deleted{\unsure{2|q}}}\in\mathbb Z$
\[
(25)\qquad\left\{\begin{aligned}
X_m^*&=\coprod_{\substack{i\in I_m\\\deleted{\unsure{d(i)}}}}X_i^*\\[1.5em]
V_{m,\deleted{\unsure{n}}\inserted{m+1}}^{\deleted{\unsure{*}}}&=\coprod_{\substack{(i,j)\in I\times I\\i\leq j,\ d(i)=m\\d(j)=m+1}}V_{i,j}\\[1.5em]
V_{m,m+1}^*&=\coprod_{\substack{i,j\text{ comme}\\\text{dessus}}}V_{ij}^*
\end{aligned}\right.
\]
de sorte qu'on a les diagrammes
\diagram{\begin{tikzcd}[column sep=large,row sep=1em]
(26) & X_m^* \arrow[r] & V_{m,m+1}^{\deleted{\unsure{*}}} & V_{m,m+1}^* \arrow[l] \arrow[dd] \\
(\mathrm{Diag})_m & & & \\
& & & X_{m+1}^*
\end{tikzcd}}
s'insérant dans un diagramme commun \unsure{\inserted{\unsure{\uline{qui n'est autre que}|qui \uline{n'est autre que}} (20)}}% doubt concerns placement after commun, not the inserted words.

infini en principe, mais que nous
écrirons sous forme finie en \deleted{\unsure{présumant|privilégiant}}
simplifier \unsure{\deleted{$I$ fini}|$I$ fini}, \inserted{$\unsure{d(I)|d(I)d}$ fini et $\subset\mathbb N$)} et en n'écrivant pas les
\deleted{\unsure{composantes} \unsure{vides}} \inserted{\unsure{partie|parties} du diagramme} \unsure{correspondant,|correspondants} à : des
sommets occupés par des espaces vides :
\margin{66}
